% =====================================================================================
%  Chapter 2 — Background and Related Work
%  Self-contained. Nothing to \input. Requires thesis/references.bib.
%
%  STRUCTURE. Written as one argument, not eighteen mini-summaries. The chain is:
%    2.1 what the policy is and why it is unsafe  ->  2.2 the control-theoretic tool
%    ->  2.3 runtime shielding for VLAs (the baseline, and its permanent cost)
%    ->  2.4 learned barriers that still optimise online
%    ->  2.5 what makes corrective data transfer
%    ->  2.6 safe RL for VLAs (scope control, not rivals)
%    ->  2.7 the exact gap
%  Do not reorder: 2.7's novelty paragraph depends on every preceding section having
%  established one of the five things it enumerates.
%
%  CITATIONS. Every claim about a paper marked [VERIFIED] in papers/FINDINGS.md was read in
%  the PDF. Entries still marked [CHECK] in references.bib need their metadata completed, but
%  the CLAIMS attached to them here are safe: they are drawn from paper_relevance_matrix.csv
%  and are positioning statements rather than numerical results.
% =====================================================================================

\chapter{Background and Related Work}\label{ch:background}

\section{Vision--Language--Action Models and the Safety Problem They Inherit}
\label{sec:bg:vla}

Vision--language--action models map camera observations and a natural-language instruction
directly to continuous robot actions, and have become the dominant paradigm for general-purpose
manipulation. The policy studied here, $\pi_{0.5}$, pairs a pretrained vision--language backbone
with a flow-matching action expert that emits chunks of continuous actions, and is trained by
heterogeneous co-training across robot demonstrations, semantic subtask prediction, and
web-derived multimodal data~\cite{physicalintelligence2025pi05}. The flow-matching action expert
is inherited from its predecessor~\cite{black2024pi0}.

What matters for this thesis is what that training does \emph{not} optimise. These models are
built and evaluated for breadth --- novel objects, unseen homes, long horizons --- and their
reported results are generalisation results. Nothing in the objective penalises contact with an
object that is not the target, and nothing in the data distribution reliably demonstrates
avoidance, because human demonstrators rarely produce near-collisions worth imitating. Broad task
competence is therefore not evidence of physical safety, and no collision-safety claim is made in
the work introducing the model~\cite{physicalintelligence2025pi05}.

The empirical consequence is stark. On the benchmark used throughout this thesis, the unmodified
policy collides with a non-target obstacle in more than eight of every ten episodes
(Section~\ref{sec:shield_efficacy}). The problem addressed here is therefore not that VLAs are bad
at manipulation; it is that competence and safety are separate properties, and only the first has
been optimised.

\section{Control Barrier Functions}\label{sec:bg:cbf}

Control theory offers a mature answer to constraint satisfaction. A control barrier function
encodes a safe set $\mathcal{C}$ as the superlevel set of a differentiable function $h$, and
converts safety into a condition that is affine in the control input: any $u$ satisfying
$\dot{h}(x,u) \ge -\alpha(h(x))$ renders $\mathcal{C}$ forward invariant, so a system that begins
safe remains safe~\cite{ames2017cbf}.

Because the condition is affine in $u$, it can be imposed as one linear inequality per constraint
inside a quadratic program that stays as close as possible to a nominal input,
\begin{equation}
  u^{\star} = \arg\min_{u} \tfrac{1}{2}\lVert u - u_{\mathrm{nom}} \rVert^{2}
  \quad \text{s.t.} \quad
  \nabla h^{\top}\!\left(f + gu\right) \ge -\alpha(h).
  \label{eq:cbf_qp}
\end{equation}
This is the object the rest of the thesis calls a \emph{shield}: a filter that accepts whatever a
nominal controller proposes and returns the minimum modification satisfying the constraint. Two
properties make it attractive as a teacher rather than merely as a guard. It is agnostic to how the
nominal action was produced, so it composes with a policy it knows nothing about; and because it
minimises $\lVert u - u_{\mathrm{nom}} \rVert$, its output is not an arbitrary safe action but
\emph{the closest safe action to what the policy wanted}, which is precisely the form a supervised
learning target should take.

Two caveats bound what may be claimed later. The guarantee is continuous-time and assumes known
control-affine dynamics, so a discrete-time implementation on a contact-rich manipulator satisfies
it only approximately (Section~\ref{sec:shield}). And forward invariance is a property of the
\emph{closed loop containing the quadratic program}. Remove the program and the guarantee is gone;
Section~\ref{sec:bg:learned} returns to how much of it can be recovered.

\section{Runtime Shielding for Vision--Language--Action Models}\label{sec:bg:shielding}

The immediate prior work applies this machinery to VLAs directly. AEGIS inserts a safety-constraint
layer between $\pi_{0.5}$ and the robot: a vision--language module grounds obstacle geometry from
observations and fused depth, and a CBF-QP modifies the policy's proposed action whenever it
violates the resulting constraint. That work also contributes SafeLIBERO, the benchmark used
here~\cite{vlsa2026}, which augments the LIBERO manipulation suite~\cite{liu2023libero} with
non-target obstacles at two difficulty levels.

The approach works, and this thesis reproduces its unshielded baseline closely --- collision
avoidance of 17.5\% against their 17.3\%, task success 58.3\% against 58.9\%, on the three suites
used here (Section~\ref{sec:shield_efficacy}). That agreement is real external validation of the
setup and is reported as such.

Two features of that work motivate everything which follows. The first is that the layer is
\emph{permanent}: it runs a quadratic program at every control step, requires obstacle geometry at
deployment, and adds a component that can itself fail. Its authors report the cost as 0.356\,ms per
step, which is modest, but the deeper cost is architectural rather than computational --- the
deployed system is a policy plus an apparatus, and the apparatus never goes away.

The second is a limitation those authors state themselves, and which this thesis measures directly.
Enforcing safety at inference can push the policy into states its training never covered, where the
policy may behave erratically and fail to recover --- an effect they name \emph{safety-induced
distribution shift}~\cite{vlsa2026}. Section~\ref{sec:res:internalise} measures exactly this as a
success cost of shielding, and it is the strongest argument that a permanent filter is not the end
state one should want. A policy that has internalised the constraint is never driven off its own
distribution by the act of being made safe.

One further difference matters for reading Chapter~\ref{ch:results}. That formulation, by its
authors' own statement, solely constrains the end-effector, and they note that unconstrained
kinematic links may consequently collide~\cite{vlsa2026}. The barrier built in
Section~\ref{sec:barrier} additionally constrains links three through seven and the hand --- an
extension whose consequences Section~\ref{sec:res:scope} quantifies.

\section{Learned Safety Representations, and Why They Still Optimise at Deployment}
\label{sec:bg:learned}

A parallel line of work removes the requirement for an analytically specified barrier by learning
one. ConBaT trains a causal transformer with a safety critic from binary safe/unsafe labels,
requiring no hand-designed constraint, and operates in embedding space rather than state space so
it applies where geometry is unavailable~\cite{meng2024conbat}. Latent Policy Barrier treats the
latent manifold of expert demonstrations as an implicit barrier, training a dynamics model on both
expert data and the policy's own rollouts, then detecting and correcting departures from that
manifold~\cite{sun2025lpb}.

Latent Policy Barrier is the closer of the two to this thesis, because its recovery data comes from
the policy's own rollout distribution --- checkpoints saved during training, executed to gather
deviations without requiring success labels, task rewards or human teleoperation. That is the same
instinct pursued here.

But both retain the property that motivated this work. ConBaT performs lightweight online
optimisation at deployment, rectifying actions by backpropagation through its critic until the CBF
conditions hold~\cite{meng2024conbat}. Latent Policy Barrier performs inference-time steering in
latent space, running gradient-based corrections through its dynamics model at every step, and
explicitly advertises plug-and-play compatibility with off-the-shelf pretrained policies without
requiring retraining or fine-tuning~\cite{sun2025lpb}. Their goal is the inverse of the one here:
they improve a frozen policy by adding machinery, where this thesis changes the policy so the
machinery can be discarded.

Between these sits the theoretical question of how much safety can transfer into weights at all.
Cosner et al.\ establish conditions under which imitating a CBF-based expert yields a learned
controller with input-to-state-safety guarantees~\cite{cosner2022il_cbf}. The result is genuine but
demanding: the expert must be a \emph{robust} CBF-QP with explicit disturbance margins, the
dynamics must be known and control-affine, and the learned controller must satisfy a
CBF-compliancy condition requiring data density along the safe-set boundary, bounded learning
error, and a known Lipschitz constant --- and even then the guarantee holds on an \emph{expanded}
safe set that degrades with data sparsity and learning error.

None of those conditions is met here, and Section~\ref{sec:disc:limitations} says which and why.
Their Remark~1 is nonetheless the cleanest available statement of this thesis's own logic: the goal
is to transfer safety guarantees from the expert controller to the learned controller rather than
to exactly mimic the expert's behaviour~\cite{cosner2022il_cbf}. Safety transfer is not imitation
fidelity --- a distinction Section~\ref{sec:res:scope} observes empirically, where the distilled
policy avoids arm-link collisions better than the teacher that supervised it.

\section{What Makes Corrective Data Transfer}\label{sec:bg:corrective}

The central question of this thesis is not whether a policy can be trained on corrected actions,
but which corrections teach. Three lines of work converge on an answer, and together they supply
the hypothesis that Chapter~\ref{ch:results} tests.

SAFE-GIL injects reachability-computed adversarial disturbances during data collection, steering a
fixed expert toward the states a learner's errors would produce so that its recovery behaviour is
recorded~\cite{ciftci2025safegil}. Their ablation is the decisive one: Gaussian noise, uniform
noise and DART all fail to reproduce the safety benefit, so it is the \emph{targeting} of the
disturbance, not augmentation as such, that matters. IntervenGen reaches the same place
differently, rolling out the current policy so that recorded mistakes are its genuine failures,
then transforming human recovery segments onto the states reached~\cite{hoque2024intervengen};
they report that ten interventions expanded synthetically outperform a hundred collected directly.

Both are worth stating carefully, because they discipline the claim this thesis makes. In SAFE-GIL
the expert is an MPC or PID controller; in IntervenGen the recovery actions come from a human
teleoperator. \textbf{Neither teacher is the learner, and in both cases safety transfers.} What the
two share is not the identity of the teacher but the location of the supervision: corrections are
supplied at states the learner itself would reach. The hypothesis this thesis inherits from them is
therefore about state coverage, and Section~\ref{sec:res:source}'s planner comparison is a test of
it rather than a demonstration that foreign controllers cannot teach.

ROAD-VLA carries the same principle into VLAs. It constructs a proximal teacher by perturbing the
student's own action-token logits with calibrated advantage estimates, and distils that teacher
along the student's on-policy rollouts, with a policy-improvement bound that holds only while the
teacher remains close to the current policy~\cite{wang2026roadvla}. Its negative result is
frequently summarised as a rejection of privileged teachers; it is not. Every rejected variant is
\emph{textual} --- a retrieved action hint, an egocentric spatial description, and a task plan ---
and the accepted teacher is derived from the student itself, so no foreign low-level controller is
tested anywhere in that work. Their explanation is that embodied post-training erodes the
backbone's in-context language ability and that a modality gap prevents discrete text from
providing the precise grounding required for continuous control, which is the reading
Section~\ref{sec:res:channels} adopts for its own language-channel negative.

Finally, one result establishes that the protocol used here is not sufficient on its own. Guerrier
et al.\ train a reinforcement-learning agent with a CBF filter active throughout training and then
remove it~\cite{guerrier2025guardrails}. Safety is not retained: the agent collides once the filter
is disabled, and their critic assigns high value \emph{inside} the obstacle. Their agent's actions
were overridden, but its objective remained a scalar reward, so no gradient ever pointed toward the
corrected action. Their curriculum variant, which blends filter and policy actions with the
filter's weight decayed to zero, does internalise the constraint.

That negative is what makes this thesis's question non-trivial. Being trained under a shield does
not make a policy safe; what the policy is trained \emph{on} decides the outcome.

\section{Policy-Level Safety Learning as an Alternative Route}\label{sec:bg:saferl}

Runtime filtering and corrective imitation are not the only options, and this thesis does not claim
they are. SafeVLA formulates VLA safety as a constrained Markov decision process, eliciting unsafe
behaviour and optimising the policy under explicit constraints~\cite{safevla}. SafeDojo performs
model-based safe reinforcement learning inside an action-conditioned video world model, estimating
imagined task progress and safety cost separately and optimising them with Lagrangian constrained
GRPO --- and evaluates on SafeLIBERO, making it the most direct methodological alternative in this
literature~\cite{safedojo}.

These matter for calibration rather than contrast. Chapter~\ref{ch:results} reports a negative
result for scalar-reward reinforcement learning in this setting; neither of these methods is
scalar-reward reinforcement learning, and both supply substantially richer constraint signals. The
negative is therefore a result about the reward design tested, and
Sections~\ref{sec:res:channels} and~\ref{sec:conc} scope it that way. The implementation of that
ablation draws on policy-gradient methods adapted to flow-matching
samplers~\cite{liu2025flowgrpo,shao2024grpo}, which belong to the methods appendix rather than to
this argument.

\section{The Gap}\label{sec:bg:gap}

The literature above establishes five things. Runtime CBF shielding makes VLAs measurably safer and
is the state of the art for this benchmark, at the cost of a permanent inference-time component and
a distribution shift its own authors document~\cite{vlsa2026}. Learned barriers relax the need for
analytic constraints but continue to optimise at
deployment~\cite{meng2024conbat,sun2025lpb}. Imitation of a CBF-safe expert can in principle
transfer formal guarantees, under conditions no large visuomotor policy currently
satisfies~\cite{cosner2022il_cbf}. Corrective data transfer when they cover the states the learner
reaches, irrespective of who produced
them~\cite{ciftci2025safegil,hoque2024intervengen,wang2026roadvla}. And exposure to a shield during
training does not, by itself, internalise anything~\cite{guerrier2025guardrails}.

What has not been asked is the amortisation question: whether the corrections a control-theoretic
shield produces can be compiled into a VLA's weights so that the shield can be \textbf{removed} at
deployment, and what bounds that transfer.

The claim of novelty must be stated precisely, because two neighbouring claims would be false.
Learning safety from CBF-related supervision is not new --- Cosner et al.\ treat it
theoretically~\cite{cosner2022il_cbf}, and learning from CBF demonstrations in simulation predates
that work. Distilling a corrective signal into a VLA is not new either; ROAD-VLA does exactly
that~\cite{wang2026roadvla}. What is new is the conjunction: a control-theoretic safety filter,
distilled from a VLA's own shielded rollouts, evaluated with the filter removed, in the setting
where runtime shielding is currently the state of the art --- together with the boundary condition
that a geometry-privileged planner, safe by construction and supplying well-formed low-level
actions, does not transfer under a matched comparison. ROAD-VLA's proximity principle predicts that
boundary; nothing in the literature tests it, because no work in this line uses a foreign low-level
controller as teacher.

Chapter~\ref{ch:method} describes the shield, the teachers, and the distillation procedure.
Chapter~\ref{ch:results} reports what transfers, what does not, and where the residual failures
live.
